\chapter{Hypokalemia (Low Potassium)}

\section*{Definition}
Hypokalemia is a lower-than-normal blood potassium concentration. Moderate or rapidly developing hypokalemia can cause muscle dysfunction and dangerous cardiac rhythm disturbances.

\section*{What Happens in Your Body}
Low extracellular potassium hyperpolarizes cell membranes and disrupts skeletal-muscle and cardiac electrical activity. It can also worsen digoxin toxicity and promote arrhythmias when magnesium is low.

\section*{Causes}
\begin{itemize}
  \item Gastrointestinal loss from vomiting, diarrhea, laxatives, or intestinal drainage.
  \item Renal loss from diuretics, mineralocorticoid excess, renal tubular disorders, or inherited salt-wasting disease.
  \item Shift into cells after insulin, beta-agonists, alkalosis, or refeeding; poor intake usually contributes rather than acts alone.
  \item Magnesium deficiency, chemotherapy, and selected antibiotics or antifungals.
\end{itemize}

\section*{When to See a Doctor}
Seek urgent care for severe weakness, paralysis, trouble breathing, fainting, palpitations, chest pain, or a markedly low result. Patients with heart disease or digoxin treatment need especially prompt assessment.

\section*{Related Symptoms}
\begin{itemize}
  \item Fatigue, muscle cramps, constipation, weakness, tingling, polyuria, and palpitations.
  \item ECG findings can include flattened T waves, ST depression, prominent U waves, and ventricular arrhythmias.
\end{itemize}

\section*{Self-Care / Home Management}
Do not self-treat a significant abnormal result with supplements without checking kidney function and the cause. Replace fluids appropriately during gastrointestinal illness and review diuretics or laxatives with a clinician.

\section*{How Your Doctor Will Evaluate You}
\begin{itemize}
  \item Repeat potassium and measure magnesium, bicarbonate, creatinine, glucose, and sometimes phosphate.
  \item Review vomiting, diarrhea, diet, medications, blood pressure, and endocrine symptoms; urine potassium helps distinguish renal from extrarenal loss.
  \item Obtain an ECG for moderate or severe hypokalemia, symptoms, rapid change, or cardiac risk.
\end{itemize}

\section*{Treatment Options}
Treatment replaces potassium orally when safe or intravenously when severe, symptomatic, or unable to take oral therapy. Magnesium deficiency must also be corrected. The cause is addressed by stopping losses, changing medications, treating diarrhea or vomiting, or managing endocrine disease.

\section*{References}
\begin{thebibliography}{99}
\bibitem{hypokalemia:ref1} Palmer BF, Clegg DJ. Physiology and pathophysiology of potassium homeostasis. \textit{Adv Physiol Educ}. 2016;40:480--490.
\bibitem{hypokalemia:ref2} Mount DB. Disorders of potassium balance. In: \textit{Harrison's Principles of Internal Medicine}. 21st ed. McGraw-Hill; 2022.
\end{thebibliography}
